%====================================================================%
% Introduction                                                       %
%                                                                    %
% 아직 산문이 없다. 아래는 **채울 자리**이고, 각 단락에 무엇을 쓸지와   %
% 어느 근거를 인용할지를 적어 두었다. 수치를 쓸 때는 보고서            %
% (실험결과_논문용_*.md)의 해당 절을 보고 그 값을 그대로 옮긴 뒤       %
% `python src/verify12.py` 로 대조할 것.                              %
%====================================================================%

\section{Introduction}

% [단락 1 — 문제] 독성 예측 모델은 보통 무작위 분할로 평가되고, 보고되는
% 지표는 거의 언제나 AUROC 다. 그런데 실제 사용은 전향적이다 — 과거 약물로
% 학습해 새 약물에 쓴다. 이 괴리를 한 문장으로 세울 것.

% [단락 2 — 알려진 것] 시간분할이 무작위 분할보다 어렵다는 것은 이미
% 알려져 있다. 선행연구를 인용하고, **그 문헌이 무엇을 보고했는지**(대개
% 판별력 저하)를 적을 것. 참고문헌은 02_참고문헌_목록.md 에 정리돼 있다.

% [단락 3 — 빈 자리] 판별력 저하만 보고하면 놓치는 것이 있다. 확률의
% 크기를 믿고 쓰는 쪽(우선순위 결정, 임계값 운용)에서는 보정이 더
% 중요한데, 보정 붕괴를 분리해 재고 그것이 유병률 변화로 설명되는지까지
% 통제한 연구가 없다는 점을 적을 것.

% [단락 4 — 이 논문이 하는 것] 두 데이터셋(DILIrank, DICTrank)에서
% 2010년 시간분할을 세우고, (i) Brier 분해로 판별력과 보정을 분리하고,
% (ii) 유병률을 맞춘 무작위 대조군 R1--R3 로 유병률 효과를 걷어내고,
% (iii) 재보정·conformal·최신성 가중 등 처방을 전향적 제약 아래에서
% 시험했다는 것을 적을 것. 숫자는 Results 로 넘기고 여기서는 설계만.

% [단락 5 — 결과 한 줄 예고] 저하의 대부분이 보정이고, 유병률을 양쪽
% 통제해도 남으며, 사후 재보정으로는 일부만 되돌아온다는 것. 구체 수치는
% 쓰지 말 것(초록과 Results 가 담당한다).

